\scanpage{5}
% One long oblique penstroke crosses the draft.
\deleted{Considérons les objets suivants

1) \quad $(P,\deleted{\unsure{\Pi|T}}\Delta)$, où $P$ est un groupe abélien libre de rang fini, $\deleted{\Pi}\Delta$ une partie finie de $P$, satisfaisant les conditions

a) \quad $\deleted{\Pi}\Delta$ engendre $P$

b) \quad Si $a\in\deleted{\Pi}\Delta$, on a $-a\in\deleted{\Pi}\Delta$, mais $na\notin\deleted{\Pi}\Delta$ si $n\ne1,-1$, $n\in\mathbb{Q}$

c) \quad Pour tout $a\in\deleted{\Pi}\Delta$, il existe une \textquotedblleft{}symétrie\textquotedblright{} $S_a$ appliquant $a$ sur $-a$, et telle que $S_a(\deleted{\Pi}\Delta)\subset\deleted{\Pi}\Delta$.

N.B. Alors les $S_a$ en question sont uniques, \deleted{\unsure{car en}}\inserted{étant en} effet \deleted{\unsure{induit}} l'unique symétrie euclidienne\inserted{ de $P^{\mathbb Q}$} pour la forme quad. suivante : $\deleted{\unsure{(x,y)=(x',y')=}}$ celle déduite par dualité de la forme $(x,y)=\sum_{a\in\Delta}\langle x,a\rangle\langle y,a\rangle$ sur $(P^{\mathbb Q})'=(P')^{\mathbb Q}$

On aura alors
\[
S_{\alpha}x=x-\langle\unsure{r|\gamma}_{\alpha},x\rangle\alpha
\]
pour un $r_{\alpha}\in P'^{\mathbb Q}$ bien déterminé. On montre que l'on a en fait $r_{\alpha}\in P'$, i.e. les $\deleted{\unsure{r}}$
\[
c_{\alpha\beta}=\unsure{-|{}}\langle r_{\alpha},\beta\rangle\qquad\text{sont }\in Z
\]
[proviendra du fait que chaque $\alpha$ est non divisible dans $P$, et ceci du fait que $\alpha$ fait partie d'un \textquotedblleft{}système fondamental de racines\textquotedblright{}]

2) \quad $(P,P',\Delta,\Delta',\varphi)$, où $P,P'$ sont deux groupes abéliens libres de rang fini \deleted{\unsure{en dualité}}\inserted{\unsure{accouplés à} $Z$}, $\Delta\subset P$, $\Delta'\subset P'$, $\varphi$ une bijection de $\Delta$ sur $\Delta'$, \deleted{\unsure{et}} les conditions \deleted{suivantes} étant vérifiées :

a) \quad $P^{\mathbb Q}$ et $P'^{\mathbb Q}$ sont en dualité, $\Delta$ engendre $P$, \inserted{$\Delta'$ engendre $P'$}

\margin{Si $\alpha'=\varphi(\alpha)$, on a $\langle\alpha',\alpha\rangle=2$}
b) \quad Pour tout $\deleted{\unsure{-}}\alpha'\in\Delta'$, \inserted{$\alpha\in\Delta$ tels que $\varphi(\alpha)=\alpha'$,} l'endomorphisme
\[
S_{\alpha}:\quad x\longmapsto x-\langle\alpha',x\rangle\alpha
\]
applique $\Delta$ dans $\Delta$, \deleted{\illegible{1}}

\noindent$\Updownarrow$\qquad b$'$) Pour tout $r\in\Delta'$, $\alpha\in\Delta$ tels que $\varphi(\alpha)=r$,
\[
S'_{\unsure{\alpha'|\alpha}}:\quad x'\longmapsto x'-\langle x',\alpha\rangle\alpha',\quad\text{invarie }\Delta'
\]}
% archive sheet number: 4
